% Packages given pdfTeX's driver option, as a document written for pdflatex
% loads them: hyperref with spaces and a braced comma in its options, and
% bookmark after it; one option list for two packages; xcolor's from
% \PassOptionsToPackage; rotating, epsfig, and colortbl, which pass it on;
% pict2e, with a slope and an arc that LaTeX's picture mode lacks; geometry;
% and \hypersetup{pdftex} after hyperref loaded. Each input's page is the
% same as without the option.
\documentclass{article}
\usepackage[pdftex]{geometry}
\PassOptionsToPackage{pdftex}{xcolor}
\usepackage[pdftex]{color,graphicx}
\usepackage[dvipsnames]{xcolor}
\usepackage[pdftex]{rotating}
\usepackage[pdftex]{epsfig}
\usepackage[pdftex]{colortbl}
\usepackage[pdftex]{pict2e}
\usepackage[ pdftex , pdftitle={A, B and C}, pdfauthor={X}, colorlinks, linkcolor=cyan]{hyperref}
\usepackage[pdftex]{bookmark}
\hypersetup{pdftex, urlcolor=blue}
\begin{document}
\section{Colors in \textcolor{ForestGreen}{green}}\label{colors}

Text in \textcolor{red}{red} and \textcolor{ForestGreen}{forest}.
A link to section~\ref{colors}, and to \href{https://example.com/}{a site}.\footnote{A
note with a link to section~\ref{colors}.}

Picture: \includegraphics[alt={wide picture}, width=2cm]{test-picture.png}.
After the picture.

Rotated: \rotatebox{90}{up}. After the rotation.

Epsfig: \epsfig{file=test-picture.png, width=1cm}. After epsfig.

\begin{tabular}{ll}
\rowcolor{yellow} one & two \\
three & four \\
\end{tabular}

\setlength{\unitlength}{1pt}
\begin{picture}(40,20)
\put(0,0){\line(7,3){30}}
\put(20,10){\arc[0,180]{10}}
\end{picture}
After the drawing.
\end{document}
